HiFi-DiffCom uses a diffusion prior conditioned on the actual channel observation,
including a confirming constraint. We retain it as an adapted generative receiver
comparison, with the forward encoder, confirming decoder and paid auxiliary
information matched to its paired SwinJSCC link. This comparison changes the
prior and recovery algorithm together; it does not isolate the superiority of
one generative backbone. The fixed16 examples are exploratory and separate from
the common500 quantitative comparison.

ARPC is closely related in its use of progressive multi-scale tokens, prediction
of unreceived scales, and autoregressive entropy probabilities. Its released
quantizer, Infinity model and text conditioning differ from our frozen system.
Accordingly, we do not claim to originate generative prefix recovery or to
outperform ARPC. Our matched internal comparisons evaluate the contribution of
permitted partial-scale transmission under the registered communication budget.
A native ARPC wireless evaluation would additionally require explicit accounting
for actual bitstreams and any receiver-side text.
